\begin{lemma}
Fix \(n\), \(\varepsilon\), and a finite base-vertex family split as
\(B_R\cup B_B\).  Put
\[
K=\noderef{TwoBiteNaturalIndependenceScale}(1+\varepsilon,n),\qquad
p=\noderef{TwoBiteEdgeProbability}\!\left(\frac12,n\right),
\]
and
\[
A=\frac{K\,n^{1/4-\varepsilon}}{6(\log n)^2}.
\]
Assume \(n>0\), \(\log n>0\), \(0<\varepsilon<1\), and
\[
|B_R|+|B_B|
\le \frac32 K n^{1/4-3\varepsilon}.
\]
Assume also the rounded-scale lower bound
\[
(1+\varepsilon)\sqrt{n\log n}\le K,
\]
the one-step mean estimate
\[
pK\le (1+\varepsilon)\log n,
\]
and the two analytic absorptions
\[
36e(1+\varepsilon)(\log n)^3\le n^{2\varepsilon},
\qquad
24(\log n)^2\le (1+\varepsilon)\sqrt{\log n}\,n^\varepsilon.
\]
Then \(A>0\), \(p\ge0\),
\[
p(|B_R|+|B_B|)K\le \frac{A}{4e},
\]
and
\[
n^{3/4-2\varepsilon}\le \frac A4.
\]
\end{lemma}
\begin{proof}
The positivity of \(A\) follows from \(K\ge(1+\varepsilon)\sqrt{n\log n}>0\),
\(n^{1/4-\varepsilon}>0\), and \(\log n>0\).  The definition of
\(\noderef{TwoBiteEdgeProbability}\) gives \(p\ge0\).

For the mean bound, the size hypothesis gives
\[
p(|B_R|+|B_B|)K
\le \frac32\,pK\,K\,n^{1/4-3\varepsilon}.
\]
Using \(pK\le(1+\varepsilon)\log n\), it is enough to absorb
\[
36e(1+\varepsilon)(\log n)^3
\]
into \(n^{2\varepsilon}\).  The stated absorption gives exactly this, and
\[
n^{1/4-3\varepsilon}n^{2\varepsilon}=n^{1/4-\varepsilon},
\]
so the result is \(p(|B_R|+|B_B|)K\le A/(4e)\).

For the lower threshold, multiply
\[
24(\log n)^2\le (1+\varepsilon)\sqrt{\log n}\,n^\varepsilon
\]
by \(n^{3/4-2\varepsilon}\).  Since
\[
\sqrt n\,n^{1/4-\varepsilon}=n^{3/4-\varepsilon}
=n^\varepsilon n^{3/4-2\varepsilon},
\]
and \(K\ge(1+\varepsilon)\sqrt{n\log n}\), this gives
\[
24(\log n)^2 n^{3/4-2\varepsilon}
\le K n^{1/4-\varepsilon}.
\]
Dividing by \(24(\log n)^2\) gives \(n^{3/4-2\varepsilon}\le A/4\).
\end{proof}
